\documentclass[11pt]{amsart}
\usepackage[localtheorems]{raeez-math-template}

\providecommand{\C}{\mathbb{C}}
\providecommand{\R}{\mathbb{R}}
\providecommand{\Z}{\mathbb{Z}}
\providecommand{\N}{\mathbb{N}}
\providecommand{\hb}{\hbar}
\providecommand{\dbar}{\bar{\partial}}
\providecommand{\g}{\mathfrak{g}}
\providecommand{\gBPS}{\mathfrak{g}_{\mathrm{BPS}}}
\providecommand{\gBPSN}{\mathfrak{g}_{\mathrm{BPS},N}}
\providecommand{\gDelta}{\mathfrak{g}_{\Delta_5}}
\providecommand{\Obs}{\mathrm{Obs}}
\providecommand{\Sym}{\mathrm{Sym}}
\providecommand{\CE}{\mathrm{CE}}
\providecommand{\Hom}{\mathrm{Hom}}
\providecommand{\chir}{\mathrm{chir}}
\providecommand{\Ran}{\mathrm{Ran}}
\providecommand{\cE}{\mathcal{E}}
\providecommand{\cF}{\mathcal{F}}
\providecommand{\cO}{\mathcal{O}}
\providecommand{\cU}{\mathcal{U}}
\providecommand{\cD}{\mathcal{D}}
\providecommand{\Aut}{\mathrm{Aut}}
\providecommand{\BV}{\mathrm{BV}}
\providecommand{\kch}{\kappa_{\mathrm{ch}}}
\providecommand{\kBKM}{\kappa_{\mathrm{BKM}}}
\providecommand{\kcat}{\kappa_{\mathrm{cat}}}
\providecommand{\kfib}{\kappa_{\mathrm{fiber}}}

\theoremstyle{plain}
\newtheorem{theorem}{Theorem}[section]
\newtheorem{proposition}[theorem]{Proposition}
\newtheorem{lemma}[theorem]{Lemma}
\newtheorem{conjecture}[theorem]{Conjecture}
\theoremstyle{definition}
\newtheorem{definition}[theorem]{Definition}
\newtheorem{construction}[theorem]{Construction}
\theoremstyle{remark}
\newtheorem{remark}[theorem]{Remark}

\title{Equivariant factorisation algebra on $K3\times E$ and the CHL
Borcherds quotient}
\author{Raeez Lorgat}
\date{2026-04-22}

\begin{document}
\maketitle

Let $S$ be a K3 surface, let $E$ be an elliptic curve, and put
$X=S\times E$. Fix holomorphic volume forms $\Omega_S$ and $dz$ and set
$\Omega_X=\Omega_S\wedge dz$. Let $\gBPS$ be a charge-graded Hall-BPS
Lie algebra for the compact $K3\times E$ branch, equipped with its
invariant pairing and with the linearisation induced by symplectic K3
automorphisms on BPS charges. For an open set $U\subset X$, write
\[
 \cE_U=\Omega_c^{0,\bullet}(U,\gBPS)[1],
 \qquad
 \cF^{\chir}_{\mathrm{BKM}}(U)
 =\CE^*_{\dbar,\chir}(\cE_U,\cO_U).
\]
The shift is the holomorphic Chern--Simons shift in complex dimension
$3$. The pairing is obtained by integrating the invariant pairing on
$\gBPS$ against $\Omega_X$.

A compact hCS comparison datum on $X$ consists of an
$\hbar$-complete nuclear completion of the Dolbeault field complex, a
Costello--Gwilliam BV quantisation solving the quantum master equation,
vanishing or cancellation of the local and global anomaly classes,
orientation data for the determinant line, a chain-level $S^3$ framing
or TCFT carrier for the compact bar operation, and a completed
hCS-to-Hall comparison preserving charge grading, pairing, coproduct,
and the relevant group linearisations. These hypotheses produce the
compact quantum object. The same formulae, without these compact
choices, define the classical local observables.

The BKM branch is the compact Hall branch. Its algebraic endpoint, once
the positive Hall half, Hopf pairing, completion, and bracket comparison
are constructed, is the Hall--Drinfeld double of the compact Hall
positive half with Borcherds denominator $\Phi_N$. The K3 self-mirror
Yangian branch is a separate Mukai-lattice branch. The
Maulik--Okounkov $E_2$ braiding enters through stable envelopes,
centres, doubles, or representation categories. The native threefold
output remains $E_1$-chiral.

\section{Total-space invariants}

\begin{proposition}[The K3$\times E$ invariant spectrum]\label{prop:k3e-four-invariants}
For $X=K3\times E$,
\[
 \kcat(X)=0,\qquad
 \kch(A_X)=0,\qquad
 \kch^{\mathrm{Heis}}(X)=3,\qquad
 \kBKM(\Phi_1)=5,\qquad
 \kfib(X;K3)=24.
\]
The numerical spectrum is $\{0,3,5,24\}$, and its entries come from
distinct constructions.
\end{proposition}

\begin{proof}
K\"unneth for coherent cohomology gives
\[
 \kcat(X)=\chi(\cO_X)
 =\chi(\cO_{K3})\chi(\cO_E)
 =(1-0+1)(1-1)=2\cdot 0=0.
\]
For the compact total-space chiral Hodge supertrace,
\[
 \kch(A_X)=\sum_{q=0}^3(-1)^q h^{0,q}(X).
\]
The Hodge polynomial in the column $H^{0,\bullet}$ is
$(1+t^2)(1+t)=1+t+t^2+t^3$, hence
\[
 \kch(A_X)=1-1+1-1=0.
\]
The Heisenberg--Cartan specialisation is rank-additive on the product
branch: the K3 factor contributes $2$ and the elliptic factor
contributes $1$, so
\[
 \kch^{\mathrm{Heis}}(X)=2+1=3.
\]
The Borcherds value is the weight of the denominator product attached
to the K3 elliptic genus. In the Gritsenko--Nikulin normalization
$c_1(0)=10$, so
\[
 \kBKM(\Phi_1)=c_1(0)/2=10/2=5.
\]
The fibre value is the Mukai rank
\[
 \kfib(X;K3)=\mathrm{rk}\,\widetilde H(K3,\Z)=24.
\]
The value $2$ is $\chi(\cO_{K3})$ on the K3 factor. It is separate from
$\kcat(X)$ and from $\kch(A_X)$.
\end{proof}

\section{Factorisation algebra}

\begin{proposition}[Holomorphic Chern--Simons factorisation]\label{prop:hcs-factorization}
The assignment $U\mapsto \cF^{\chir}_{\mathrm{BKM}}(U)$ is the
classical holomorphic Chern--Simons factorisation algebra on $X$. With
the compact hCS comparison datum it is the corresponding quantum compact
factorisation algebra.
\end{proposition}

\begin{proof}
The fields $U\mapsto\cE_U$ form the elliptic local moduli problem for
holomorphic Chern--Simons theory with coefficient Lie algebra $\gBPS$.
The differential is $\dbar$, the bracket is the pointwise bracket on
$\gBPS$ combined with wedge product of $(0,\bullet)$-forms, and the
local functional is
\[
 I_U(\alpha)=\int_U \Omega_X\wedge
 \left(\frac{1}{2}\langle \alpha,\dbar\alpha\rangle
 +\frac{1}{6}\langle \alpha,[\alpha,\alpha]\rangle\right).
\]
All operations are local on $U$. For disjoint opens
$U_1,\ldots,U_m\subset V$, extension by zero and multiplication of
chiral Chevalley--Eilenberg cochains give the factorisation maps
\[
 \bigotimes_{i=1}^m
 \CE^*_{\dbar,\chir}(\cE_{U_i},\cO_{U_i})
 \longrightarrow
 \CE^*_{\dbar,\chir}(\cE_V,\cO_V).
\]
Associativity follows from associativity of wedge product and the
Chevalley--Eilenberg coalgebra. Descent is Costello--Gwilliam descent
for observables of an elliptic local moduli problem. On polydiscs this
is the standard holomorphic Chern--Simons factorisation algebra. On
overlaps the transition functions preserve $\dbar$, the bracket, the
pairing, $\Omega_X$, and the chosen compact BV data, so the local
objects glue. The compact zero-mode, anomaly, orientation, framing,
completion, and hCS-to-Hall comparison hypotheses are exactly the
additional data needed for the quantum compact object.
\end{proof}

\section{Equivariance}

\begin{definition}[$G$-equivariant factorisation algebra]\label{def:g-equivariant}
Let a group $G$ act on $X$ by holomorphic maps preserving $\Omega_X$ and
on $\gBPS$ by pairing-preserving Lie algebra automorphisms. A
$G$-equivariant factorisation algebra on $X$ is a factorisation algebra
$\cF$ together with quasi-isomorphisms
\[
 \rho_g:g^{-1}\cF\xrightarrow{\sim}\cF,\qquad g\in G,
\]
compatible with the factorisation maps and satisfying
$\rho_{gh}=\rho_g\circ g^{-1}\rho_h$.
\end{definition}

\begin{proposition}[Volume-preserving equivariance]\label{prop:volume-equivariance}
Let $G\subset\Aut(X,\Omega_X)$ be a finite subgroup equipped with a
pairing-preserving linearisation on $\gBPS$. Assume that $G$ preserves
the completion, counterterms, orientation datum, framing or TCFT carrier,
and hCS-to-Hall comparison map. Then
$\cF^{\chir}_{\mathrm{BKM}}$ has a canonical $G$-equivariant structure.
\end{proposition}

\begin{proof}
For $g\in G$, pullback gives a cochain isomorphism
\[
 g^*:\Omega_c^{0,\bullet}(U,\gBPS)
 \xrightarrow{\sim}
 \Omega_c^{0,\bullet}(g^{-1}U,\gBPS).
\]
Here $g^*$ means pullback on forms together with the given automorphism
of $\gBPS$. Because $g$ is holomorphic, $g^*$ commutes with $\dbar$.
Because $g$ preserves $\Omega_X$ and the pairing on $\gBPS$, it
preserves the BV pairing and the local holomorphic Chern--Simons
functional. Preservation of the compact datum identifies the chosen
orientation, counterterms, framing, and completion with their pullbacks.
Therefore $g^*$ induces
quasi-isomorphisms on chiral Chevalley--Eilenberg observables. The
identity $(gh)^*=h^*g^*$ gives the group law, and functoriality
of pullback gives compatibility with disjoint-union factorisation
maps.
\end{proof}

\section{The CHL quotient}

Fix $N\in\{1,2,3,4,6\}$. Let $g_N\in\Aut_{\mathrm{s}}(K3)$ be a finite
symplectic automorphism of order $N$, and let $t_N\in E[N]$ have exact
order $N$. Define
\[
 \gamma_N(s,e)=(g_Ns,e+t_N),
 \qquad
 \Gamma_N=\langle\gamma_N\rangle\cong\Z/N\Z .
\]
The action is taken with the induced $\Gamma_N$-linearisation on
$\gBPS$. Let $\gBPSN=(\pi_{N,*}\underline{\gBPS})^{\Gamma_N}$ denote
the descended locally constant Lie algebra on the quotient.

\begin{lemma}[Free diagonal action]\label{lem:free-diagonal-action}
The action of $\Gamma_N$ on $X=K3\times E$ is free, and
\[
 X_N=X/\Gamma_N
\]
is a smooth compact threefold with trivial canonical bundle. Its
holomorphic volume form is the descent of $\Omega_X$.
\end{lemma}

\begin{proof}
If $\gamma_N^m(s,e)=(s,e)$ for $0<m<N$, then $e+mt_N=e$, hence
$mt_N=0$. Since $t_N$ has exact order $N$, this is impossible. Thus no
nontrivial element of $\Gamma_N$ has a fixed point. The quotient is
therefore smooth and the projection $\pi_N:X\to X_N$ is finite
\'etale. The automorphism $g_N$ is symplectic, so $g_N^*\Omega_{K3}
=\Omega_{K3}$; translation preserves $dz$. Hence
$\gamma_N^*\Omega_X=\Omega_X$, and $\Omega_X$ descends to a nowhere
vanishing holomorphic $3$-form on $X_N$. This is the Calabi--Yau
condition used here.
\end{proof}

\begin{definition}[CHL chiral CE complex]\label{def:chl-ce}
For an open set $V\subset X_N$, define
\[
 \cF^{\chir}_{\mathrm{BKM}}(X_N)(V)
 =
 \CE^*_{\dbar,\chir}
 \bigl(\Omega_c^{0,\bullet}(V,\gBPSN)[1],\cO_V\bigr).
\]
Equivalently,
\[
 \cF^{\chir}_{\mathrm{BKM}}(X_N)(V)
 \cong
 \left(
 \CE^*_{\dbar,\chir}
 \bigl(\Omega_c^{0,\bullet}(\pi_N^{-1}V,\gBPS)[1],
 \cO_{\pi_N^{-1}V}\bigr)
 \right)^{\Gamma_N}.
\]
The second description is the $g_N$-twined model. The freeness of the
action makes it quotient descent and eliminates fixed-locus inertia
sectors.
\end{definition}

\begin{proposition}[Factorisation descent to the CHL quotient]\label{prop:chl-descent}
There is an isomorphism of factorisation algebras on $X_N$:
\[
 \cF^{\chir}_{\mathrm{BKM}}(X_N)
 \cong
 \left(\pi_{N,*}\cF^{\chir}_{\mathrm{BKM}}(X)\right)^{\Gamma_N}.
\]
\end{proposition}

\begin{proof}
For every $V\subset X_N$, finite \'etale descent identifies
\[
 \Omega_c^{0,\bullet}(V,\gBPSN)
 \cong
 \Omega_c^{0,\bullet}(\pi_N^{-1}V,\gBPS)^{\Gamma_N}.
\]
The group is finite and the ground field has characteristic zero, so
taking $\Gamma_N$-invariants is exact. The chosen nuclear completion
makes the completed chiral Chevalley--Eilenberg functor commute with
these finite invariants and with its defining projective limits. The
identification respects $\dbar$, the pointwise $\gBPS$ bracket, the
pairing, and the factorisation maps. The displayed isomorphism follows
open by open and is compatible with disjoint unions.
\end{proof}

\begin{proposition}[Compact invariants of the CHL quotient]\label{prop:chl-hodge-invariants}
For every $N\in\{1,2,3,4,6\}$,
\[
 \kcat(X_N)=0,\qquad
 \kch(A_{X_N})=0.
\]
\end{proposition}

\begin{proof}
Since $\pi_N:X\to X_N$ is finite etale,
\[
 H^q(X_N,\cO_{X_N})
 \cong H^q(X,\cO_X)^{\Gamma_N}.
\]
The action on $H^{0,\bullet}(K3)$ is trivial in degrees $0$ and $2$:
$g_N$ is symplectic and fixes $\Omega_{K3}$. Translation on $E$ is
homotopic to the identity and acts trivially on $H^{0,0}(E)$ and
$H^{0,1}(E)$. Thus the invariant Hodge column is still
\[
 h^{0,0}=h^{0,1}=h^{0,2}=h^{0,3}=1.
\]
Therefore
\[
 \kch(A_{X_N})=\sum_q(-1)^q h^{0,q}(X_N)=1-1+1-1=0,
 \qquad
 \kcat(X_N)=\chi(\cO_{X_N})=0.
\]
\end{proof}

\section{Borcherds denominator}

Let
\[
 \phi_N(\tau,z)=\phi_{0,1}^{(g_N)}(\tau,z)
 =\sum_{n,\ell}c_N(n,\ell)q^n y^\ell
\]
be the $g_N$-twined K3 elliptic genus in the Gritsenko--Nikulin CHL
normalization, and write $c_N(0)=c_N(0,0)$. The Borcherds product
$\Phi_N=\mathrm{Borch}(\phi_N)$ has weight $c_N(0)/2$. In the five CHL
orders used here,
\[
\begin{array}{c|ccccc}
N&1&2&3&4&6\\ \hline
c_N(0)&10&8&6&4&2\\
\kBKM(\Phi_N)=c_N(0)/2&5&4&3&2&1 .
\end{array}
\]

\begin{proposition}[BKM weight]\label{prop:bkm-weight}
For $N\in\{1,2,3,4,6\}$,
\[
 \kBKM(\Phi_N)=\mathrm{wt}(\Phi_N)=\frac{c_N(0)}{2}.
\]
The root multiplicities of the Borcherds--Kac--Moody algebra attached
to $\Phi_N$ are the Fourier coefficients of $\phi_N$ appearing as
exponents in the Borcherds product.
\end{proposition}

\begin{proof}
Borcherds' singular theta correspondence gives the product expansion
of $\Phi_N$ and identifies its weight with one half of the constant
coefficient of the vector-valued input at the chosen cusp. Gritsenko
and Nikulin apply this construction to the $g_N$-twined K3 elliptic
genus for the CHL orders $N\in\{1,2,3,4,6\}$. The coefficient table
above is the corresponding constant-term table. By definition
$\kBKM$ is the automorphic weight of this denominator product, so
the displayed equality follows. The same product formula writes each
positive root factor with exponent given by a Fourier coefficient
$c_N(n,\ell)$ of $\phi_N$, after the usual Weyl-chamber choice.
\end{proof}

\begin{remark}[Borcherds invariant formula]\label{rem:borcherds-invariant-formula}
The Borcherds weight is read from the denominator product constant term.
For $N=1$,
\[
 \kBKM(\Phi_1)=5,\qquad
 \kch(A_{K3\times E})=0,\qquad
 \chi(\cO_E)=0.
\]
The Heisenberg value $\kch^{\mathrm{Heis}}(K3\times E)=3$ and the
K3 value $\chi(\cO_{K3})=2$ belong to different constructions. The
Borcherds side is governed by the single identity
$\kBKM(\Phi_N)=c_N(0)/2$.
\end{remark}

\section{Chiral BKM comparison}

\begin{proposition}[Output scope at $d=3$]\label{prop:d3-output-scope}
The factorisation algebra on $X_N$ is an $E_1$-chiral object. Any
$E_2$-structure belongs to a center, Drinfeld double, module-category
enhancement, or representation-category construction after the required
hypotheses are supplied; it is outside the native $d=3$ output.
\end{proposition}

\begin{proof}
The CY-to-chiral functor is dimension-stratified: $d\leq 2$ gives
$E_2$-chiral output, while $d\geq 3$ gives $E_1$-chiral output. The
quotient $X_N$ has complex dimension $3$ and trivial canonical bundle by
Lemma~\ref{lem:free-diagonal-action}; hence the native output is $E_1$.
Passing from an $E_1$-algebra to an $E_2$-object requires a separate
center, double, or module-enhancement construction and the functorial
hypotheses for that construction.
\end{proof}

\begin{conjecture}[Chiral envelope of the CHL BKM algebra]\label{conj:chl-chiral-bkm}
For $N\in\{1,2,3,4,6\}$, the $E_1$-chiral factorisation algebra
$\cF^{\chir}_{\mathrm{BKM}}(X_N)$ admits a comparison map to the compact
Hall positive half. After the Hall--Drinfeld double, completion,
nondegenerate Hopf pairing, and vacuum-envelope data are supplied, this
comparison identifies the completed target with the chiral envelope of
the Borcherds--Kac--Moody algebra whose denominator is $\Phi_N$. The
denominator root multiplicities are those of
Proposition~\ref{prop:bkm-weight}.
\end{conjecture}

The required hypotheses are the equivariant $d=3$ CY-to-chiral
construction on $X_N$ with compact framing and TCFT data, the orientation
and completion package for the hCS-to-Hall comparison, construction of
the compact Hall positive half with nondegenerate Hopf pairing and
completed double, the comparison between the reduced DT/BPS Hall algebra
of $X_N$ and the Borcherds product $\Phi_N$, and functorial preservation
of the $\Gamma_N$-linearisation through the chiral envelope.

\paragraph{References.}
Costello--Gwilliam, \emph{Factorization Algebras in Quantum Field
Theory}, Vols.~I--II, for descent of observables of elliptic local
moduli problems. Nikulin 1980, \emph{Izv.\ Akad.\ Nauk SSSR Ser.\ Mat.}
\textbf{44}, for finite symplectic automorphisms of K3 surfaces. Mukai
1988, \emph{Invent.\ Math.}~\textbf{94}, for finite symplectic K3
automorphism groups and Mathieu embeddings. Chaudhuri--Hockney--Lykken
1995, \emph{Phys.\ Rev.\ Lett.}~\textbf{75}, for CHL quotients.
Borcherds 1998, \emph{Invent.\ Math.}~\textbf{132}, and
Gritsenko--Nikulin, \emph{Automorphic forms and Lorentzian
Kac--Moody algebras}, for the weight formula
$\mathrm{wt}(\Phi_N)=c_N(0)/2$. Cheng--Duncan--Harvey 2014,
\emph{Res.\ Math.\ Sci.}~\textbf{1}, for twined K3 elliptic genera.

\end{document}
